\section{Introdução}

As métricas DORA descrevem velocidade e estabilidade da entrega de software
\cite{forsgren2018accelerate,dora}. Em projetos open-source, porém, uma release não
comprova implantação em produção e uma falha de CI não comprova incidente. O objetivo
deste estudo é medir proxies explícitos a partir do GitHub e avaliar quanto os
resultados dependem da definição operacional, sem interpretar associação como causa.
Investigamos sete questões e um bônus (Tabela~\ref{tab:questoes}).

\begin{table}[htbp]
\centering\footnotesize
\caption{Questões e expectativas registradas antes da análise (H1--H7).}
\label{tab:questoes}
\begin{tabular}{p{.07\textwidth}p{.43\textwidth}p{.40\textwidth}}
\hline
RQ & Questão & Hipótese informal \\
\hline
01 & Qual a frequência das entregas? & H1: mediana menor que uma release/semana, com cauda longa. \\
02 & Qual o tempo entre commit e entrega? & H2: lead time por release bem maior que por commit. \\
03 & Qual a taxa de falha? & H3: CFR de CI maior que CFR por releases corretivas. \\
04 & Qual o tempo de recuperação? & H4: poucas horas, cauda pesada e censura não desprezível. \\
05 & Frequência associa-se a falhas? & H5: associação fraca e negativa. \\
06 & Quais características se associam às métricas? & H6: tipo e contribuidores mais associados que linguagem. \\
07 & A classificação depende da definição? & H7: concordância moderada/alta, com trocas próximas aos cortes. \\
08 & Qual a proporção de releases corretivas? & Bônus descritivo: rework rate, sem nova hipótese prévia. \\
\hline
\end{tabular}
\end{table}

A H2 supõe commits concentrados perto da publicação. A H3 antecipa que CI bloqueie
problemas antes de uma entrega. H5 acompanha a ideia de que velocidade e estabilidade
não são necessariamente um trade-off~\cite{doraMetricas}; H6 considera a capacidade
de manutenção das comunidades. H7 antecipa que proxies diferentes possam alterar as
categorias. As formulações completas anteriores aos resultados permanecem no
histórico do repositório; não foram reescritas para concordar com os achados.
